\documentclass[11pt,fleqn,oneside,a4paper]{article}
\usepackage[T1]{fontenc}
\usepackage[utf8]{inputenc}
\usepackage{color}
\usepackage[ngerman]{babel}
\usepackage{textcomp}
\usepackage{fancyhdr}
\pagestyle{fancy}

% page layout
\setlength{\paperwidth}{210mm}
\setlength{\paperheight}{297mm}
\setlength{\voffset}{0mm}
\setlength{\hoffset}{0mm}
\setlength{\oddsidemargin}{0mm}
\setlength{\evensidemargin}{0mm}
\setlength{\marginparsep}{1.5mm}
\setlength{\marginparwidth}{60mm}
\setlength{\marginparpush}{1mm}
\setlength{\parindent}{0cm} %% avoid indent on new paragraph
\setlength{\oddsidemargin}{-17.9mm}
\setlength{\headsep}{3mm}
\setlength{\headheight}{8mm}
\setlength{\topskip}{0mm}
\setlength{\topmargin}{-18mm}
\setlength{\footskip}{-1mm}
\setlength{\headwidth}{189mm}
\setlength{\textwidth}{189mm}
\setlength{\textheight}{261.5mm}
\renewcommand{\baselinestretch}{1}

% rulewidth
\renewcommand{\headrulewidth}{0.25mm}
\renewcommand{\footrulewidth}{0.25mm}
\setlength{\fboxrule}{0.1mm}

\lhead{\bf Zugangsdaten}
\chead{\LARGE \bf {{ schoolclass }}}
\rhead{\bf \large \today}
\lfoot{ {{ school_longname }} }
\cfoot{}
\rfoot{ {{ filename }} }

\begin{document}

{% for data in datablocks -%}
    {% set block_loop = loop %}
    {% for student in data -%}

Guten Tag {{ student['firstname'] }} {{ student['lastname'] }},

\vspace{2mm}

Der/Die Netzwerkbetreuer/in

{{ admins_print }}

heißen Sie an der Schule \textbf{ {{ school_longname }} } willkommen.

\vspace{2mm}

Hier sind Ihre Zugangsdaten zum pädagogischen Netz.

Bitte ändern Sie Ihr Passwort möglichst bald in der Schulkonsole (siehe unten).

\begin{center}
\fbox{
\begin{minipage}[b][53,98mm][c]{85,6mm} % scheckkarte
\begin{center}
\Large {{ student['lastname'] }}, {{ student['firstname'] }} \par Klasse: {{ student['schoolclass'] }} \par Login: {{ student['login'] }} \par Erst-Passwort: {{ student['password'] }} \par Schule: {{ school_longname }}
\end{center}
\end{minipage}
}
\end{center}

\kern 3mm \hrule width \hsize height 0.25mm \kern 0.5mm \hrule width \hsize height 0.25mm \kern 3mm
Anmeldung an einem Rechner im pädagogischen Netz:\\[-7mm]

\begin{tabbing}
\qquad Benutzername:\qquad \= \texttt{ {{ student['login'] }} } \\
\qquad Kennwort:     \> \texttt{ {{ student['password'] }} } \\
\qquad Domäne:       \> \texttt{ {{ domain }} } (falls erforderlich)\\
\end{tabbing}
\vspace{-9mm}
\kern 3mm \hrule width \hsize height 0.25mm \kern 0.5mm \hrule width \hsize height 0.25mm \kern 3mm

Startseite (Landing Page) für die Schule:\\[-7mm]
\begin{tabbing}
\qquad URL:\qquad\qquad \= \texttt{ {{ urlstart_print }} } \\
\qquad {{ urlstart_comment_print }} \\
\end{tabbing}
\vspace{-9mm}
\kern 3mm \hrule width \hsize height 0.25mm \kern 0.5mm \hrule width \hsize height 0.25mm \kern 3mm

\medskip

\rule{\linewidth}{0.25mm}
Pädagogische Funktionen mit der Schulkonsole im Intranet:\\[-7mm]
\begin{tabbing}
\qquad URL:\qquad\qquad \= \texttt{ {{ urlschuko_print }} } \\
\qquad {{ urlschuko_comment_print }} \\
\end{tabbing}
\vspace{-9mm}
\rule{\linewidth}{0.25mm}

Abruf der schulischen E-Mail:\\[-7mm]
\begin{tabbing}
\qquad URL:\qquad\qquad \= \texttt{ {{ urlmail_print }} } \\
\qquad {{ urlmail_comment_print }} \\
\end{tabbing}
\vspace{-9mm}
\rule{\linewidth}{0.25mm}

Lernplattform mit Moodle:\\[-7mm]
\begin{tabbing}
\qquad URL:\qquad\qquad \= \texttt{ {{ urlmoodle_print }} } \\
\qquad {{ urlmoodle_comment_print }} \\
\end{tabbing}
\vspace{-9mm}
\rule{\linewidth}{0.25mm}

    {% endfor -%}
    {% if not block_loop.last -%}
        \newpage
    {% endif %}
{% endfor %}
\end{document}
